% Volume V: Admissibility and Continuation
% Part X. Public Institutions as Epistemic Ecologies
% Chapter 56: Archives and Institutional Memory
% Spine: proof-spines/ch056-spine.md

\chapter{Archives and Institutional Memory}
\label{ch:ch056}

Archives preserve material beyond its present salience. Their democratic function is to keep future reconstruction possible, so the chapter models archival selection, retention, indexing, and degradation as active determinants of what a polity can later know.

An archive is not a warehouse but a future-facing admissibility structure:
\[
\mathcal A_t
\longrightarrow
\{\text{questions answerable at }t+\Delta\}.
\]
What is kept, discarded, misfiled, or allowed to decay changes the set of questions that later investigators, courts, journalists, or citizens can responsibly answer. Institutional memory is therefore a designed capacity, not a passive remainder.
